% Section 3, part 6 of 8: the gauge identity (G), the canonical gauge, symplecticity, HPF.

\begin{thm}[The gauge identity (G)]
\label{thm:gau-G}
For every $P\in\Gau(A,\tilde A)$, in particular for every regular gauge from the pulled-back model to the quintic system,
\begin{equation}
\label{eq:gau-G}
P(u,z)\,\YC(v(u),z)=P_{11}(0,z)\,\YCQ(u,z)\qquad\text{in }\bbQ\lser z\fps u^5 .
\end{equation}
\begin{enumerate}
\item[(a)] For a graded regular gauge $P$, $P_{11}(0,z)$ is a unit of $\bbQ\fps{z^5}$, $\Pc:=P\,P_{11}(0,z)^{-1}$ is a graded regular gauge, and $\Pc$ is the unique regular gauge $P'$ with $P'\,\YC(v(u),z)=\YCQ$.
\item[(b)] For a regular gauge $P$, $P\,\YC(v(u),z)$ has no positive power of $z$ if and only if $P_{11}(0,z)$ is a constant; the regular gauges with this property are the $\mu\Pc$, $\mu\in\bbQ^\times$.
\item[(c)] For any two cyclic systems $A,\tilde A$ with one simple pole over $\bbQ\fps u$ and every $P\in\Gau(A,\tilde A)$: $P\,Y^{\mathrm C}_{\tilde A}=P_{11}(0,z)\,Y^{\mathrm C}_A$.
\end{enumerate}
\end{thm}

\begin{proof}
We prove (c), of which \eqref{eq:gau-G} is the case of the two systems of this subsection: $\YC(v(u),z)=Y^{\mathrm C}_{\tilde A}$ by \cref{lem:gau-uniq}, both being solutions of $z\de_uY=\tilde AY$ with value $-\mathbf e_1$ (\cref{lem:gau-vanmod,lem:gau-vanq}~(c)). Put $\tilde Y:=Y^{\mathrm C}_{\tilde A}\in\bbQ[z^{-1}]\fps u^5$. The product $P\tilde Y$ lies in $\bbQ\lser z\fps u^5$: at each order in $u$ it is a finite sum of products of a power series in $z$ and a polynomial in $z^{-1}$. It solves the first system: $z\de_u(P\tilde Y)=(AP-P\tilde A)\tilde Y+P\tilde A\tilde Y=A(P\tilde Y)$. Its value at $u=0$ is $P(0,z)(-\mathbf e_1)=-P_{11}(0,z)\mathbf e_1$ by \cref{lem:gau-R1}. $P_{11}(0,z)Y^{\mathrm C}_A$ is also a solution, $P_{11}(0,z)$ being constant in $u$, with the same value by \cref{lem:gau-vanq}~(c). \Cref{lem:gau-uniq} with $B=A$ gives the identity.

(a) $P_{11}(0,z)$ is a unit of $\bbQ\fps{z^5}$ by \cref{prop:gau-graded}~(iii), so $\Pc$ is a graded regular gauge by \cref{prop:gau-graded}~(ii), and $(\Pc)_{11}(0,z)=1$; by \eqref{eq:gau-G}, $\Pc\YC(v(u))=\YCQ$. Conversely, let $P'$ be any regular gauge with $P'\YC(v(u))=\YCQ$. By \cref{thm:gau-U}~(b), (d), $P'=\Pc\cn(z)$ for some $\cn\in\bbQ\fps z^\times$, so that $\YCQ=\cn(z)\Pc\YC(v(u))=\cn(z)\YCQ$, i.e.\ $(\cn(z)-1)\YCQ=0$ in $\bbQ\lser z\fps u^5$; the coefficient of $u^0$ is $-(\cn(z)-1)\mathbf e_1$, so $\cn=1$. (b) If $P_{11}(0,z)=\mu$ is constant, then $P\YC(v(u))=\mu\YCQ$ lies in $\bbQ[z^{-1}]\fps u^5$. Conversely, the coefficient of $u^0$ of $P\YC(v(u))$ is $-P_{11}(0,z)\mathbf e_1$ with $P_{11}(0,z)\in\bbQ\fps z$, so it has no positive power of $z$ only if $P_{11}(0,z)$ is constant. With $P=\Pc\cn(z)$ we have $P_{11}(0,z)=\cn(z)$, which gives the last assertion.
\end{proof}

\begin{defn}[The canonical gauge]
\label{def:gau-can}
In a conifold frame, the \emph{canonical gauge} $\Pc$ is the regular gauge from the pulled-back model to the quintic system with $(\Pc)_{11}(0,z)=1$. It exists and is unique: $\Pc=P\,P_{11}(0,z)^{-1}$ for any graded regular gauge $P$ (\cref{thm:gau-G}~(a)), and a regular gauge $\Pc\cn(z)$ (\cref{thm:gau-U}~(b)) has $(\Pc\cn)_{11}(0,z)=\cn(z)$. Equivalently, $\Pc$ is the unique regular gauge with $\Pc\YC(v(u))=\YCQ$. It is graded, rational (\cref{thm:gau-U}~(d)), colour covariant, $\operatorname{Rot}_i\Pc=\Pc$ (\cref{prop:gau-graded}~(iv)), and $p_{0,a}(0)=1$ for its diagonal $z^0$-term (\cref{lem:gau-diag}). We write $\Pc=\sum_kz^k\sum_bp_{k,b}(u)E_{b+k,b}$.
\end{defn}

\begin{prop}[Symplecticity]
\label{prop:gau-symp}
Let $A,\tilde A$ be self-dual cyclic systems with one simple pole over $\Bbbk\fps u$ with (EP) and $r=\tilde r$, and let $P$ be a regular gauge from $\tilde A$ to $A$, $\cn(z):=P_{11}(0,z)$. Then
\begin{equation}
\label{eq:gau-symp}
P(u,-z)^T\eta\,P(u,z)=\cn(z)\cn(-z)\,\eta .
\end{equation}
In particular, for the two systems of \cref{sec:gau-van}, $\Pc$ is symplectic, $\Pc(u,-z)^T\eta\Pc(u,z)=\eta$; the symplectic regular gauges are the $\Pc\cn(z)$ with $\cn(z)\cn(-z)=1$, so symplecticity constrains only the even part of $\log(\cn/\cn(0))$; for graded gauges, $\cn=\pm\exp(a_1z^5+a_3z^{15}+a_5z^{25}+\cdots)$. Among the symplectic regular gauges, only (G) singles out $\Pc$.
\end{prop}

\begin{proof}
For $S:=P(u,-z)^T\eta P(u,z)$, the gauge equation at $-z$, transposed, and self-duality $A^T\eta=\eta A$, $\tilde A^T\eta=\eta\tilde A$ (\cref{def:gau-onepole}) give
\[
zS'=\bigl(-P(-z)^TA^T+\tilde A^TP(-z)^T\bigr)\eta P(z)+P(-z)^T\eta\bigl(AP(z)-P(z)\tilde A\bigr)=\tilde A^TS-S\tilde A,
\]
so $T:=\eta^{-1}S$ satisfies $zT'=\tilde AT-T\tilde A$ and $T=t(z)\,1$ by \cref{thm:gau-U}~(b) for $(\tilde A,\tilde A)$. At $u=0$, by \cref{lem:gau-R1} (column $1$ of $P(0,z)$ is $\cn(z)\mathbf e_1$, and $P_{22}(0,\cdot)=\cn$ as $r=\tilde r$) and $\eta_{a1}=5\delta_{a2}$, $S_{21}(0,z)=5P_{22}(0,-z)\cn(z)=5\cn(-z)\cn(z)$, while $S_{21}=5t(z)$. The consequences follow from $(\Pc\cn)_{11}(0,z)=\cn(z)$: $\cn(z)\cn(-z)=1$ forces $\cn(0)=\pm1$ and is equivalent to the oddness of $\log(\cn/\cn(0))$, and graded gauges have $\cn\in\Bbbk\fps{z^5}$.
\end{proof}

%-------------------------------------------------------------------------------
\subsection{The normalisation series}
\label{sec:gau-R0}

We compare $\Pc$ applied to the model columns with the rows of \cref{def:set-rows}: colour $0$ in \cref{thm:gau-R0}, the others by colour covariance (\cref{cor:gau-colours}). Here $\theta_u:=u\,\dd/\dd u$, $\theta_{\zh}:=\zh\,\dd/\dd\zh$, and a series $\Phi(\zh,u)=\sum_n\zh^n\mathrm{ph}_n(u)$ with $\mathrm{ph}_n\in\bbQ\lser s$ is identified with $\sum_nz^ns^{-4n}\mathrm{ph}_n$ through $\zh=z/s^4=L^4z$ (\cref{def:set-points}); $\Dq$ acts on $\bbQ\lser s$ as in \cref{lem:gau-confexp}.

\begin{lem}[HPF]
\label{lem:gau-HPF}
For $e\in\bbZ$ let $\mathscr O_e:=1+\zh(1-u)\bigl(\tfrac45\theta_{\zh}-\tfrac e5-\theta_u\bigr)$, and for $\Phi=\sum_n\zh^n\mathrm{ph}_n$ as above put $\Phi^{[0]}:=\Phi$ and $\Phi^{[j]}:=\mathscr O_{-j}\Phi^{[j-1]}$, $j=1,\dots,5$.
\begin{enumerate}
\item $L\Phi$ satisfies $\PFfive^{(0)}(L\Phi)=0$, where $\PFfive^{(0)}$ is $\PFfive$ with $z\Dq$ replaced by $z\Dq+L$ (that is, $\re^{u_0/z}L\Phi$ solves $\PFfive$ formally), if and only if
\begin{equation}
\label{eq:gau-HPF}
\Phi^{[5]}-\Phi-(1-u)\Bigl[3\zh\Phi^{[4]}+\tfrac{17}5u\zh^2\Phi^{[3]}+\tfrac95u^2\zh^3\Phi^{[2]}+\tfrac{274}{625}u^3\zh^4\Phi^{[1]}+\tfrac{24}{625}u^4\zh^5\Phi\Bigr]=0 .
\end{equation}
\item The coefficient of $\zh^{n+1}$ of the left side of \eqref{eq:gau-HPF} is $(1-u)\bigl[(4n-5\theta_u)\mathrm{ph}_n+r_n\bigr]$, where, with $\Phi^\flat_{j,n}$ the coefficient of $\zh^n$ of $\Phi^{[j]}$ computed with $\mathrm{ph}_n$ replaced by $0$ (and $\Phi^{[j]}_m:=\coef{\zh^m}\Phi^{[j]}$, zero for $m<0$),
\[
r_n:=\sum_{j=1}^5\Bigl(\frac{4n+j}5-\theta_u\Bigr)\Phi^\flat_{j-1,n}-3\Phi^\flat_{4,n}-\tfrac{17}5u\Phi^{[3]}_{n-1}-\tfrac95u^2\Phi^{[2]}_{n-2}-\tfrac{274}{625}u^3\Phi^{[1]}_{n-3}-\tfrac{24}{625}u^4\mathrm{ph}_{n-4}
\]
depends only on $\mathrm{ph}_0,\dots,\mathrm{ph}_{n-1}$. So \eqref{eq:gau-HPF} holds if and only if $(4n-5\theta_u)\mathrm{ph}_n=-r_n$ for all $n\geq0$; here $r_0=0$, and $r_1=\frac35$ when $\mathrm{ph}_0=1$.
\item If every $\mathrm{ph}_n$ lies in $\bbQ\fps u$, \eqref{eq:gau-HPF} reads $(4n-5k)\coef{u^k}\mathrm{ph}_n=-\coef{u^k}r_n$ for all $n,k\geq0$: for $5k\neq4n$ the coefficient $\coef{u^k}\mathrm{ph}_n$ is fixed by $\mathrm{ph}_0,\dots,\mathrm{ph}_{n-1}$, while for $n=5m$ the coefficient $\fc_m:=\coef{u^{4m}}\mathrm{ph}_{5m}$ is free and $\coef{u^{4m}}r_{5m}=0$ is required. In particular $\mathrm{ph}_0$ is a constant, and $\mathrm{ph}_0=1$ forces $\mathrm{ph}_1=-\frac3{20}$.
\end{enumerate}
\end{lem}

\begin{proof}
(1) From $\Dq(s^cf)=(1-u)s^{c-5}\bigl(-\frac c5-\theta_u\bigr)f$ (as $\Dq=-(1-u)u^{-1}\theta_u$), the coefficient of $z^m$ of $(z\Dq+L)(s^eF^\flat)$, $F^\flat=\sum_m\zh^mF^\flat_m$ (a local letter), is $s^{e-4m-1}\bigl[F^\flat_m+(1-u)\bigl(\frac{4(m-1)-e}5-\theta_u\bigr)F^\flat_{m-1}\bigr]$; so $(z\Dq+L)(s^eF^\flat)=s^{e-1}\mathscr O_eF^\flat$. (2) By induction $(z\Dq+L)^j(s^{-1}\Phi)=s^{-1-j}\Phi^{[j]}$; as $z=s^4\zh$, $z^m(z\Dq+L)^{5-m}(L\Phi)$ is $s^{-1}u^{m-1}\zh^m\Phi^{[5-m]}$. (3) With $M=z\Dq$, $5^5q=1-u$ and the elementary symmetric functions $(\sigma_1,\dots,\sigma_5)=(3,\frac{17}5,\frac95,\frac{274}{625},\frac{24}{625})$ of $\frac15,\dots,\frac55$, $\PFfive=uM^5-1-(1-u)\sum_{k=1}^5\sigma_kz^kM^{5-k}$; replacing $M$ by $z\Dq+L$, applying to $L\Phi$ and using (2) gives $s^{-1}$ times the left side of \eqref{eq:gau-HPF}, which is (i).
(4) By definition $\Phi^{[j]}_m=\Phi^{[j-1]}_m+(1-u)\bigl(\frac{4(m-1)+j}5-\theta_u\bigr)\Phi^{[j-1]}_{m-1}$, so $\Phi^{[5]}_{n+1}-\mathrm{ph}_{n+1}=(1-u)\sum_{j=1}^5(\frac{4n+j}5-\theta_u)\Phi^{[j-1]}_n$; with $\Phi^{[j]}_n=\mathrm{ph}_n+\Phi^\flat_{j,n}$ the terms in $\mathrm{ph}_n$ are $(1-u)\bigl[\sum_{j=1}^5(\frac{4n+j}5-\theta_u)-3\bigr]\mathrm{ph}_n=(1-u)(4n-5\theta_u)\mathrm{ph}_n$, and the rest is $(1-u)r_n$. Clearly $r_0=0$; for $\mathrm{ph}_0=1$, $\Phi^\flat_{j,1}=(1-u)\frac{j(j+1)}{10}$ and $(c-\theta_u)(1-u)=c(1-u)+u$ give $r_1=\sum_{j=2}^5\frac{j(j-1)}{10}\bigl(\frac{4+j}5(1-u)+u\bigr)-6(1-u)-\frac{17}5u=\frac35$.
(5) (iii) follows from $\theta_uu^k=ku^k$: for $n=0$ it gives $\mathrm{ph}_0\in\bbQ$, and for $n=1$, with $r_1=\frac35$, $\mathrm{ph}_1=-\frac3{20}$.
\end{proof}

% Section 3, part 7 of 8: polynomiality, formal WKB uniqueness, colour covariance of the model columns and of the actual rows.

\begin{lem}[Power-series solutions of HPF are polynomial]
\label{lem:gau-poly}
Let $\Phi$ be a solution of \eqref{eq:gau-HPF} with $\mathrm{ph}_0=1$ and every $\mathrm{ph}_n\in\bbQ\fps u$.
\begin{enumerate}
\item Every $\mathrm{ph}_n$ is a polynomial in $u$.
\item For $\mathsf k\in\bbQ\fps z$ with $\mathsf k(0)=1$, the product $\Phi\,\mathsf k(z)$ (with $z^n=s^{4n}\zh^n$) solves \eqref{eq:gau-HPF}, and its coefficients lie in $\bbQ\fps u$ if and only if $\mathsf k\in\bbQ\fps{z^5}$.
\item The solutions of \eqref{eq:gau-HPF} with $\mathrm{ph}_0=1$ and coefficients in $\bbQ\fps u$ are exactly the $\Phi\,\mathsf k(z)$, $\mathsf k\in1+z^5\bbQ\fps{z^5}$; for every choice of the free coefficients $(\fc_m)_{m\geq1}$ there is exactly one of them, and $\Phi^\circ$ denotes the one with all $\fc_m=0$.
\end{enumerate}
In particular, as soon as one solution with coefficients in $\bbQ\fps u$ and $\mathrm{ph}_0=1$ exists, the resonance conditions $\coef{u^{4m}}r_{5m}=0$ hold for every choice of the free coefficients.
\end{lem}

\begin{proof}
(i) Induction on $n$: $r_n$ is obtained from polynomials by $\theta_u$ and multiplication by polynomials, and by \cref{lem:gau-HPF}~(iii) $\coef{u^k}\mathrm{ph}_n=-\coef{u^k}r_n/(4n-5k)$ vanishes for $k>\deg r_n$, except possibly the free coefficient at $k=4n/5$. (ii) $\theta_{\zh}z^m=mz^m$ and $\theta_uz^m=\frac{4m}5z^m$, so every $\mathscr O_e$, hence the left side of \eqref{eq:gau-HPF}, commutes with multiplication by a series in $z$. If $n_0$ is the least index with $5\nmid n_0$ and $\mathsf k_{n_0}\neq0$, the coefficient of $\zh^{n_0}$ of $\Phi\mathsf k$ contains the fractional power $\mathsf k_{n_0}s^{4n_0}$ plus power series. (iii) For a second such solution $\Phi'$, suppose $\Phi'\equiv\Phi\mathsf k^{(N)}$ modulo $\zh^N$ with $\mathsf k^{(N)}\in1+z^5\bbQ[z^5]$ ($\mathsf k^{(1)}=1$). Both sides solve \eqref{eq:gau-HPF} by (ii) and have the same $r_N$, so the difference $\delta$ of their coefficients of $\zh^N$ satisfies $(4N-5\theta_u)\delta=0$: $\delta=0$ if $5\nmid N$, and $\delta=\delta_0u^{4m}$ if $N=5m$, absorbed by $\mathsf k^{(N+1)}:=\mathsf k^{(N)}+\delta_0z^{5m}$, as $z^{5m}\Phi=u^{4m}\zh^{5m}(1+O(\zh))$. The limit satisfies $\Phi'=\Phi\mathsf k$, uniquely as $\Phi$ is a unit; the converse is (ii). The free coefficients $\fc_m(\Phi\mathsf k)=\mathsf k_{5m}+\sum_{0<m'<m}\mathsf k_{5m'}\fc_{m-m'}(\Phi)+\fc_m(\Phi)$ form a triangular system with unit diagonal in the $\mathsf k_{5m}$, which gives the rest.
\end{proof}

\begin{lem}[Formal WKB uniqueness]
\label{lem:gau-wkb}
Let $\Bbbk\supseteq\bbQ(\zeta)$ be a field (for $i=0$, any field $\Bbbk\supseteq\bbQ$), and let $(\cR,\Dq,L)$ be $(\Bbbk\fps q,\,q\,\dd/\dd q,\,(1-5^5q)^{-1/5})$ or $(\Bbbk\lser s,\Dq,s^{-1})$; in both, $\Dq L=(L^6-L)/5$ and $5^5q=1-L^{-5}$. For $i\in\bbZ/5$ let $\PFfive^{(i)}\in\cR[z][\Dq]$ be $\PFfive$ with $z\Dq$ replaced by $z\Dq+\zeta^iL$; a \emph{colour-$i$ WKB amplitude} of $\PFfive$ over $\cR$ is a series $\cA=\sum_{n\geq0}\cA_nz^n\in\cR\fps z$ with $\PFfive^{(i)}\cA=0$ (so that $\re^{u_i/z}\cA$ solves $\PFfive$ formally, $\Dq u_i=\zeta^iL$).
\begin{enumerate}
\item $\PFfive^{(i)}=\sum_{k=1}^5z^k\mathcal T^{\mathrm w}_k$ with $\mathcal T^{\mathrm w}_k\in\cR[\Dq]$ (local letters), and $\mathcal T^{\mathrm w}_1(Lg)=5\zeta^{4i}\Dq g$ for $g\in\cR$.
\item If $\cA,\cA'$ are colour-$i$ WKB amplitudes with $\cA_0=\cA'_0=L$, then $\cA'=\cA\,\mathsf k(z)$ for a unique $\mathsf k\in1+z\Bbbk\fps z$; for $i=0$ and $\cA,\cA'$ with coefficients in $\bbQ\fps q$ or $\bbQ\lser s$, $\mathsf k\in1+z\bbQ\fps z$.
\item $\cA(z)\mapsto\cA(\zeta^{-i}z)$ is a bijection from colour-$0$ to colour-$i$ WKB amplitudes that preserves $\cA_0$.
\end{enumerate}
\end{lem}

\begin{proof}
(i) The coefficient of $z^0$ is $(\zeta^iL)^5-1-q(5\zeta^iL)^5=L^5(1-5^5q)-1=0$. As $\Dq L=(L^6-L)/5$, the coefficients of $z$ of $(z\Dq+\zeta^iL)^5$ and of $\prod_{k=1}^5\bigl(5\zeta^iL+z(5\Dq+k)\bigr)$ act by
\[
f\mapsto\zeta^{4i}\bigl(5L^4\Dq f+2(L^9-L^4)f\bigr),\qquad f\mapsto5^4\zeta^{4i}(25L^4\Dq f+10L^9f+5L^4f).
\]
With $5^4q=(1-L^{-5})/5$, $\mathcal T^{\mathrm w}_1f=\zeta^{4i}\bigl(5L^{-1}\Dq f+(L^{-1}-L^4)f\bigr)$, and $\mathcal T^{\mathrm w}_1(Lg)=5\zeta^{4i}\Dq g$.
(ii) At $z^{n+1}$ the equation is $\mathcal T^{\mathrm w}_1\cA_n=-\sum_{k=2}^5\mathcal T^{\mathrm w}_k\cA_{n+1-k}$, and $\ker\mathcal T^{\mathrm w}_1=\Bbbk L$ by (i), the constants of $\Dq$ on $\cR$ being $\Bbbk$. Multiplication by a constant series commutes with $\PFfive^{(i)}$. If $\cA'\equiv\cA\mathsf k^{(n)}$ modulo $z^n$ with $\mathsf k^{(n)}(0)=1$, then $B:=\cA'-\cA\mathsf k^{(n)}$ is an amplitude vanishing modulo $z^n$, so $\mathcal T^{\mathrm w}_1B_n=0$, $B_n=\mathsf k_nL$, and $\mathsf k^{(n+1)}:=\mathsf k^{(n)}+\mathsf k_nz^n$ works. The limit is unique since $\cA$ is a unit; for rational data of colour $0$ every $\mathsf k_n$ is rational.
(iii) $\sigma_i(\cA)(z):=\cA(\zeta^{-i}z)$ is a ring automorphism of $\cR\fps z$ commuting with $\Dq$, with $\sigma_i\circ(z\Dq+L)=\zeta^{-i}(z\Dq+\zeta^iL)\circ\sigma_i$; apart from its constant term $\PFfive=M^5-1-q\prod_k(5M+kz)$, $M=z\Dq$, is homogeneous of degree $5$ in $(M,z)$, so $\sigma_i\circ\PFfive^{(0)}=\PFfive^{(i)}\circ\sigma_i$.
\end{proof}

\begin{lem}[Colour covariance of the model columns]
\label{lem:gau-colmod}
For $i\in\bbZ/5$ let $F^{(i)}=(F^{(i)}_b)_{b\in\bbZ/5}$ be the solution with $F^{(i)}_0(0)=1$ of the colour-$i$ T-system of \cref{prop:set-modrows}~(b), with $\mathcal T^{(i)}_k=-\zeta^i+\frac t5(k-4t\,\dd/\dd t)$; thus $\Fs_b=F^{(0)}_b$.
\begin{enumerate}
\item $F^{(i)}_b(t)=\zeta^{ib}\Fs_b(\zeta^{-i}t)$.
\item The rows $\Rhat^{\mathrm m,(i)}_b(t)=\zeta^{i(b+1)}\Fs_b(\zeta^{-i}t)=\zeta^iF^{(i)}_b(t)$ of \cref{def:set-M0} are the rows of a WKB solution of the model system with exponential $\re^{U_i/z}$, $U_i:=-\frac54\zeta^i\varsigma^4$: the \emph{model columns} $\Qr^{\mathrm m,(i)}_b(v,z):=\varsigma^{\nu_b-1}\Rhat^{\mathrm m,(i)}_b(z/\varsigma^4)$ satisfy $z\,\de_v\Qr^{\mathrm m,(i)}_b+(\de_vU_i)\Qr^{\mathrm m,(i)}_b=b^{\mathrm m}_b\Qr^{\mathrm m,(i)}_{b+1}$.
\item At $v=v(u)$: $\Qr^{\mathrm m,(i)}_b(u,z)=\zeta^{i(b+1)}\Qr^{\mathrm m,(0)}_b(u,\zeta^{-i}z)$, i.e.\ $\Qr^{\mathrm m,(i)}=\operatorname{Rot}_i\Qr^{\mathrm m,(0)}$.
\end{enumerate}
\end{lem}

\begin{proof}
(i) As $t\,\dd/\dd t$ commutes with $t\mapsto\zeta^{-i}t$, $\mathcal T^{(i)}_k[G(\zeta^{-i}t)]=\zeta^i[\mathcal T^{(0)}_kG](\zeta^{-i}t)$. Hence $G_b(t):=\zeta^{ib}\Fs_b(\zeta^{-i}t)$ satisfies the five relations of the colour-$i$ T-system (the closing one using $\zeta^{-5i}=1$) and $G_0(0)=1$, so $F^{(i)}=G$ by the uniqueness in \cref{prop:set-modrows}~(b) (pivot $4n\zeta^{4i}\neq0$). (ii) By \cref{prop:set-modrows}~(b), $\re^{U_i/z}\varsigma^{\nu_b-1}F^{(i)}_b(z/\varsigma^4)$ solves the model system, and so does its multiple by $\zeta^i$. (iii) $\varsigma$ depends on $u$ only.
\end{proof}

\begin{lem}[Colour covariance of the actual rows]
\label{lem:gau-colrows}
For every $i\in\bbZ/5$, as identities of formal series in $z$ with coefficients in $\bbQ(\zeta)\fps q$,
\[
\Phi_i(\zh,u)=\Phi_0(\zeta^{-i}\zh,u)\qquad\text{and}\qquad \Qr^{(i)}_a(z)=\zeta^{i(a+1)}\Qr^{(0)}_a(\zeta^{-i}z),\ \text{i.e.}\ \Qr^{(i)}_{a,n}=\zeta^{i(a+1-n)}\Qr^{(0)}_{a,n},
\]
where $\Qr^{(i)}_{a,n}:=\coef{z^n}\Qr^{(i)}_a$ are the coefficients of the rows of \cref{def:set-rows}.
\end{lem}

\begin{proof}
\emph{Step 1.} By \cref{cor:set-cyclicPF5}~(b), $L\Phi_i$ is a colour-$i$ WKB amplitude in the sense of \cref{lem:gau-wkb}, with coefficient $L$ at $z^0$, since $R=1+O(z)$ gives $\coef{z^0}\Qr^{(i)}_0=\eta(1,\tilde e_i)=\zeta^iL/C_0$ (\cref{cor:set-cyclicPF5}~(a), \cref{def:set-frob}). By \cref{lem:gau-wkb}~(iii), so is $L\Phi_0(\zeta^{-i}\zh,u)$ ($\zh=L^4z$). \emph{Step 2.} By \cref{lem:gau-wkb}~(ii), $L\Phi_0(\zeta^{-i}\zh,u)=L\Phi_i(\zh,u)\,\mathsf k(z)$ with $\mathsf k\in1+z\bbQ(\zeta)\fps z$. \emph{Step 3.} At $q=0$, where $u=L=C_0=1$ and $\zh=z$, (QRR) (\cref{thm:set-qrr}) gives $\Phi_i(z,1)=h_{\bw_i}(z)$ and $\Phi_0(\zeta^{-i}z,1)=h_{\bw_0}(\zeta^{-i}z)=h_{\bw_i}(z)$ (\cref{lem:set-hrot}); as $h_{\bw_i}$ is a unit and $\mathsf k$ does not depend on $q$, $\mathsf k=1$. This step is the only place in this section where (QRR) enters. \emph{Step 4.} For $X^{(i)}_a:=\zeta^{-i}\Qr^{(i)}_a$, $X^{(i)}_0=L\Phi_i/C_0$ (\cref{def:set-rows}), so $X^{(i)}_0(z)=X^{(0)}_0(\zeta^{-i}z)$ by Step 3. The forward relations of \cref{prop:set-cyclicLR} read $X^{(i)}_{a+1}=(C_1\kappa^{\mathrm{LR}}_a)^{-1}(z\Dq+\zeta^iL)X^{(i)}_a$ with coefficients independent of $z$ and $i$, and $(z\Dq+\zeta^iL)[F(\zeta^{-i}z)]=\zeta^i[(z\Dq+L)F](\zeta^{-i}z)$; by induction on $a$, $X^{(i)}_a(z)=\zeta^{ia}X^{(0)}_a(\zeta^{-i}z)$. \emph{Step 5.} Hence $\Qr^{(i)}_a(z)=\zeta^{i(a+1)}\Qr^{(0)}_a(\zeta^{-i}z)$; take the coefficient of $z^n$.
\end{proof}

% Section 3, part 8 of 8: the normalisation series D, the rows in the ring, the colours of D, (F_g).

\begin{thm}[The normalisation series $\Dn$; colour $0$]
\label{thm:gau-R0}
Let $(\alpha,\beta)$ be a conifold frame, $\Pc$ the canonical gauge, and $\Qr^{\mathrm m,(i)}$ the model columns of \cref{lem:gau-colmod} at $v=v(u)$.
\begin{enumerate}
\item For every $i$, $\Pc\Qr^{\mathrm m,(i)}$ is a formal WKB solution of the quintic system with exponential $\re^{u_i/z}$: its rows $X_a\in\bbQ(\zeta)\lser s\fps z$ satisfy $(z\Dq+\zeta^iL)X_a=\kappa_aX_{a+1}$ for $a\in\bbZ/5$, with $L=s^{-1}$. With $\Lambda_s:=\diag(s^{1-\nu_a})$, $\Lambda_\varsigma:=\diag(\varsigma^{1-\nu_b})$ and $\Ghat(u,\zh):=\Lambda_s\Pc(u,s^4\zh)\Lambda_\varsigma^{-1}$ (the hatted gauge of \cref{lem:mai-Ghat}),
\[
\Lambda_s\,\Pc\,\Qr^{\mathrm m}\big|_{z=s^4\zh}=\Ghat(u,\zh)\,\Rhat^{\mathrm m}(\tau(u)\zh),\qquad \coef{\zh^k}\Ghat_{b+k,b}=\tau^{(1-\nu_b)/4}u^{\mu(b,k)}p_{k,b}(u),
\]
with $\mu(b,k)=(4k+\nu_b-\nu_{b+k})/5\in\bbZ$, $\mu(b,k)\geq\max(0,(4k-4)/5)$; this is a power series in $u$ whose coefficients are power series in $\zh$.
\item $\Phican(\zh,u):=\varpi_0\cdot\bigl[\Ghat(u,\zh)\Rhat^{\mathrm m}(\tau\zh)\bigr]_{00}$ (row $0$, colour $0$) solves \eqref{eq:gau-HPF}, with coefficients in $\bbQ\fps u$ and $\mathrm{ph}^{\mathrm{can}}_0=1$. Hence every $\mathrm{ph}^{\mathrm{can}}_n$ is a polynomial in $u$ (\cref{lem:gau-poly}) and $A^{\mathrm{can}}_n:=L^{4n+1}\mathrm{ph}^{\mathrm{can}}_n(L^{-5})$ lies in $\bbQ[L,L^{-1}]$.
\item Let $A^{\mathrm{act}}:=C_0\Qr^{(0)}_0=\sum_nA^{\mathrm{act}}_nz^n$ be the flat amplitude of the actual colour-$0$ row $0$ (\cref{def:set-rows}), with $A^{\mathrm{act}}_n\in\bbQ(\zeta)\fps q$. In $\bbQ(\zeta)\fps q$, with $L=(1-5^5q)^{-1/5}$, both $A^{\mathrm{can}}:=\sum_nA^{\mathrm{can}}_nz^n$ and $A^{\mathrm{act}}$ are colour-$0$ WKB amplitudes of $\PFfive$ with coefficient $L$ at $z^0$. Hence $A^{\mathrm{can}}=A^{\mathrm{act}}\,\Dn_0(z)$ for a unique $\Dn_0\in1+z\bbQ(\zeta)\fps z$, and $A^{\mathrm{act}}_n\in\bbQ(\zeta)[L,L^{-1}]$ for every $n$.
\item In $\bbQ(\zeta)[L,L^{-1}]\fps z$, embedded in $\bbQ(\zeta)\lser s\fps z$ by $L\mapsto s^{-1}$, the identity of (iii) holds; in hatted form it reads $\Phiact(\zh,u)\,\Dn_0(s^4\zh)=\Phican(\zh,u)$, $\Phiact:=\Phi_0=L^{-1}A^{\mathrm{act}}$. The rows of the actual colour-$0$ column have conifold expansions (\cref{cor:gau-rowsNN}), they satisfy $(z\Dq+L)X_a=\kappa_aX_{a+1}$, and
\begin{equation}
\label{eq:gau-R0zero}
\Qr^{(0)}\,\Dn_0(z)=\Pc\,\Qr^{\mathrm m,(0)}(v(u)),
\end{equation}
with $\Dn_0$ constant in $u$ and the same exponential $\re^{u_0/z}$ stripped on both sides.
\end{enumerate}
\end{thm}

\begin{proof}
(i) By \cref{lem:gau-colmod}~(ii) and the chain rule, $z\de_u\Qr^{\mathrm m,(i)}+(\de_uU_i)\Qr^{\mathrm m,(i)}=\tilde A\Qr^{\mathrm m,(i)}$ with $U_i=U_i(v(u))=-\frac54\zeta^iu^{4/5}/\tau=u_i-u_i^{\mathrm c}$ (\cref{def:set-tau}, \cref{lem:set-canon}), so $\de_uU_i=\de_uu_i$. For $X:=\Pc\Qr^{\mathrm m,(i)}$ the gauge equation gives $z\de_uX+(\de_uu_i)X=AX$, which is $(z\Dq+\zeta^iL)X_a=\kappa_aX_{a+1}$ after multiplication by $-(1-u)$; the amplitudes of $\re^{u_i/z}$ and $\re^{U_i/z}$ agree, as $u_i-U_i=u_i^{\mathrm c}$ is constant. Hatting: $\varsigma=s\tau^{-1/4}$ (\cref{def:set-tau}), so $z/\varsigma^4=\tau\zh$, $\Lambda_\varsigma\Qr^{\mathrm m}=\Rhat^{\mathrm m}(\tau\zh)$, and the entry $(b+k,b)$ of $\Ghat$ at $\zh^k$ is $s^{1-\nu_{b+k}}s^{4k}p_{k,b}\,s^{\nu_b-1}\tau^{(1-\nu_b)/4}$. Since $\nu_a\equiv-a$ and $\nu\in\{-1,\dots,3\}$, $\mu(b,k)\in\bbZ$, $\mu(b,k)\geq(4k-4)/5$ and $\mu(b,0)=0$.

(ii) As $\nu_0=0$, $L\Phican=\varpi_0(\Pc\Qr^{\mathrm m,(0)})_0$. By (i) and \cref{cor:gau-cycsys}~(ii) ($\cR=\bbQ\lser s\lser z$, $w=L$, $(y_0,y_1)=(\varpi_0,\varpi_1)$), this solves $\PFfive^{(0)}$, so $\Phican$ solves \eqref{eq:gau-HPF} by \cref{lem:gau-HPF}~(i). Colour $0$ is rational ($\Rhat^{\mathrm m,(0)}_b=\Fs_b$, \cref{def:set-M0}; \cref{def:gau-can}), so by (i) the coefficients lie in $\bbQ\fps u$; $\Ghat(u,0)$ is diagonal with $(0,0)$ entry $\tau^{1/4}p_{0,0}=1/\varpi_0$ (\cref{lem:gau-diag}) and $\Fs_0(0)=1$, so $\mathrm{ph}^{\mathrm{can}}_0=1$. \Cref{lem:gau-poly}~(i) applies, and $A^{\mathrm{can}}_n=u^{-(4n+1)/5}\mathrm{ph}^{\mathrm{can}}_n(u)\in\bbQ[L,L^{-1}]$.

(iii) $L\mapsto(1-5^5q)^{-1/5}$ and $L\mapsto s^{-1}$ are injective homomorphisms of differential rings on $\bbQ(\zeta)[L,L^{-1}]$, $\Dq L=(L^6-L)/5$ (the first as $L-1$ has order $1$ in $q$). The colour-$0$ WKB equations of $A^{\mathrm{can}}$ are identities in $\bbQ[L,L^{-1}]$ that hold after $L\mapsto s^{-1}$ by (ii) ($\sum_nA^{\mathrm{can}}_nz^n=L\Phican$), hence in $\bbQ\fps q$; and $A^{\mathrm{can}}_0=L$. $A^{\mathrm{act}}$ is a colour-$0$ WKB amplitude by \cref{cor:set-cyclicPF5}~(b), with $A^{\mathrm{act}}_0=C_0\eta(1,\tilde e_0)=L$ (\cref{def:set-frob}). \Cref{lem:gau-wkb}~(ii) with $\Bbbk=\bbQ(\zeta)$ gives $\Dn_0$, and $A^{\mathrm{act}}_n=\sum_{m\leq n}A^{\mathrm{can}}_m\coef{z^{n-m}}\Dn_0^{-1}$ is the image of a unique Laurent polynomial.

(iv) By injectivity the identity of (iii) holds in $\bbQ(\zeta)[L,L^{-1}]\fps z$, hence in $\bbQ(\zeta)\lser s\fps z$; divide by $L$ and put $z=s^4\zh$. Row $0$ of the actual colour-$0$ column is represented by $A^{\mathrm{act}}_n(L)c_0^{-1}$, and rows $a\geq1$ by the forward relations of \cref{prop:set-cyclicLR}, with coefficients represented in $\Nring$; this is \cref{cor:gau-rowsNN}~(i) for colour $0$, and by \cref{prop:gau-quintic}~(i) the expansions satisfy $(z\Dq+L)X_a=\kappa_aX_{a+1}$ for $a\leq3$. Row $0$ of \eqref{eq:gau-R0zero}: $A^{\mathrm{act}}\Dn_0/\varpi_0=A^{\mathrm{can}}/\varpi_0=L\Phican/\varpi_0=(\Pc\Qr^{\mathrm m,(0)})_0$ by (ii). Rows $1,\dots,4$ of both sides follow from row $0$ by $X_{a+1}=\kappa_a^{-1}(z\Dq+L)X_a$ ($\Dn_0$ is constant; (i)), as does the closing relation on the left.
\end{proof}

\begin{cor}[The rows lie in the differential ring]
\label{cor:gau-rowsNN}
\begin{enumerate}
\item For all $i,a,n$ the row coefficient $\Qr^{(i)}_{a,n}$ is the evaluation (\cref{def:set-confexp}) of an explicit $\logq$-free element of $\Nring\otimes\bbQ(\zeta)$: for colour $0$, row $0$ is $A^{\mathrm{act}}_n(L)\,c_0^{-1}$ with $A^{\mathrm{act}}_n\in\bbQ(\zeta)[L,L^{-1}]$ (\cref{thm:gau-R0}~(iii)); the rows $a=1,\dots,4$ follow from $\Qr_{a+1,n}=(C_1\kappa^{\mathrm{LR}}_a)^{-1}(\Dq\Qr_{a,n-1}+L\Qr_{a,n})$, with $(C_1\kappa^{\mathrm{LR}}_a)^{-1}$ represented by $(c_0^2/\mathsf W,\ \mathsf W^2L^{-5}/c_0^2,\ c_0^2/\mathsf W,\ c_0^{-1})$; and colour $i$ is $\zeta^{i(a+1-n)}$ times colour $0$ (\cref{lem:gau-colrows}). These representatives are used throughout; in this sense the rows lie in $\Nring\otimes\bbQ(\zeta)$ and have conifold expansions.
\item Consequently every polynomial over $\bbQ(\zeta)$ in finitely many row coefficients and in the $\Delta_i^{\pm1/2}=(\zeta^{4i}C_0L^{-1})^{\pm1}$ has a $\logq$-free representative; this applies to the edges, leaves and genus-$g$ graph terms built from the rows (\cref{def:mai-Fcf}).
\item The covariance $\Qr^{(i)}_{a,n}=\zeta^{i(a+1-n)}\Qr^{(0)}_{a,n}$ holds for the representatives, hence for the conifold expansions and for the hatted rows $\hat{\Qr}_a=s^{1-\nu_a}\Qr_a$, $z=s^4\zh$.
\end{enumerate}
\end{cor}

\begin{proof}
(i) The $\logq$-free subring of $\Nring$ (generated by $L^{\pm1}$, the $c_j$, the $i_j$, $c_0^{-1}$, $\mathsf W^{-1}$) is $\Dq$-stable: $\PFfour c_0=0$ does not involve $\logq$, and in $\Conc(c_0,c_0\logq+i_0)=0$, which defines $\Dq i_2$, the coefficient of $\logq$ is $\sum_{a,b}c_a\Jet_{ab}c_b=0$ by antisymmetry of $\Jet$ (\cref{lem:gau-conc}~(iii)). The forward relations at $z^n$ give the recursion, with $\logq$-free coefficients (\cref{prop:gau-quintic}~(i)), and $\ev$ is a homomorphism of differential rings; colour $i$ is \cref{lem:gau-colrows}. (ii) $\Delta_i^{1/2}$ is represented by the unit $\zeta^{4i}c_0L^{-1}$. (iii) The conifold expansion is a $\bbQ(\zeta)$-linear ring homomorphism, and hatting multiplies $\coef{z^n}$ of row $a$ by $s^{1-\nu_a+4n}$ in every colour.
\end{proof}

\begin{cor}[The colours of $\Dn$]
\label{cor:gau-colours}
For every $i\in\bbZ/5$, $\Dn_i(z):=\Dn_0(\zeta^{-i}z)$ satisfies, in $\bbQ(\zeta)\lser s\fps z$ and for the conifold expansions of the rows,
\[
\Qr^{(i)}\,\Dn_i(z)=\Pc\,\Qr^{\mathrm m,(i)}(v(u)),
\]
and it is the unique series in $\bbQ(\zeta)\fps z$ with this property. Hence, with $\Dn:=\diag(\Dn_0,\dots,\Dn_4)$ and the matrices of rows (columns indexed by the colours),
\begin{equation}
\label{eq:gau-R0}
\Qr\,\Dn=\Pc\,\Qr^{\mathrm m}(v(u)).
\end{equation}
Every $\Dn_i$ lies in $1+z\bbQ(\zeta)\fps z$, with $\coef{z^n}\Dn_i=\zeta^{-in}\coef{z^n}\Dn_0$; so $\Dn_0=1$ implies $\Dn=1$. In particular the expanded rows of every colour form a WKB solution of the quintic system with exponential $\re^{u_i/z}$.
\end{cor}

\begin{proof}
By \cref{lem:gau-colmod}~(iii), $\operatorname{Rot}_i\Pc=\Pc$ (\cref{def:gau-can}), the product rules of \cref{prop:gau-graded}~(iv), \eqref{eq:gau-R0zero} and \cref{cor:gau-rowsNN}~(iii),
\[
\Pc\Qr^{\mathrm m,(i)}=\operatorname{Rot}_i\bigl(\Pc\Qr^{\mathrm m,(0)}\bigr)=\operatorname{Rot}_i\bigl(\Qr^{(0)}\Dn_0(z)\bigr)=\operatorname{Rot}_i(\Qr^{(0)})\,\Dn_0(\zeta^{-i}z)=\Qr^{(i)}\,\Dn_0(\zeta^{-i}z).
\]
Uniqueness: $\Qr^{(i)}_0$ is not a zero divisor, its coefficient of $z^0$ being $\zeta^is^{-1}/\varpi_0\neq0$. The coefficients of $\Dn_i$ follow from \cref{thm:gau-R0}~(iii), and $\Qr^{(i)}=\Pc\Qr^{\mathrm m,(i)}\Dn_i^{-1}$ is a WKB solution by \cref{thm:gau-R0}~(i).
\end{proof}

\begin{defn}[The window $\Fg{g}$]
\label{def:gau-Fg}
For $g\geq2$, $\Fg{g}$ is the statement $\Dn_0\equiv1\pmod{z^{3g-2}}$; by \cref{cor:gau-colours} it implies $\Dn\equiv1\pmod{z^{3g-2}}$. It is the hypothesis of the comparison theorem (\cref{thm:flh-main}). By \cref{thm:gam-D}, $\Dn=1$, so $\Fg{g}$ holds for every $g\geq2$.
\end{defn}


